\documentclass[11pt]{article}
\usepackage[margin=1.04in,top=.94in,bottom=.97in]{geometry}
\usepackage{amsmath,amssymb,amsthm,booktabs,mathtools,microtype}
\usepackage[dvipsnames]{xcolor}
\definecolor{inkblue}{HTML}{194A69}
\usepackage[colorlinks=true,linkcolor=inkblue,citecolor=inkblue,urlcolor=inkblue]{hyperref}
\hypersetup{pdftitle={Identifying abelian number fields from a short interval of zeta zeros},pdfauthor={Research note prepared for Vladimir Reshetnikov}}
\usepackage{fancyhdr}
\pagestyle{fancy}\fancyhf{}\fancyhead[L]{\small\itshape A short interval of zeta zeros}\fancyhead[R]{\small Dirichlet \(L\)-functions}\fancyfoot[C]{\thepage}
\setlength{\headheight}{14pt}\renewcommand{\headrulewidth}{.25pt}
\newtheorem{theorem}{Theorem}[section]
\newtheorem{lemma}[theorem]{Lemma}
\newtheorem{corollary}[theorem]{Corollary}
\newcommand{\C}{\mathbb C}\newcommand{\R}{\mathbb R}\newcommand{\Q}{\mathbb Q}
\newcommand{\dd}{\,d}\newcommand{\supp}{\operatorname{supp}}
\newcommand{\sinc}{\operatorname{sinc}}\newcommand{\ord}{\operatorname{ord}}
\newcommand{\one}{\mathbf1}
\DeclareMathOperator{\sech}{sech}
\setlength{\emergencystretch}{2em}
\begin{document}
\thispagestyle{plain}
\noindent{\small\sffamily\color{inkblue} ANALYTIC NUMBER THEORY\hfill 27 NOVEMBER 2026}
\vspace{12pt}

\noindent{\LARGE\bfseries Identifying abelian number fields from a short interval of zeta zeros\par}
\vspace{5pt}
\noindent{\large A conductor-uniform separation theorem for Dirichlet \(L\)-functions}
\vspace{12pt}

\noindent\textbf{Abstract.}
For each fixed \(d\), an abelian extension of \(\Q\) of degree at most \(d\) and conductor at most \(Q\) is determined by the ordinates, with multiplicities, of its Dedekind zeta zeros in
\(0<\gamma<C_d(\log Q)^2/\log\log Q\) for all sufficiently large \(Q\).
More generally, on every interval of this length plus an additive triple logarithm of its height, the zero multiplicity vectors of any fixed number of distinct primitive Dirichlet characters have full complex rank. We prove a uniform lower bound on every linear combination and a stability bound for imperfectly matched lists. Dirichlet GRH enters through the location of zeros and a smoothed prime sum. The gain in the conductor aspect comes from averaging interference from nearby prime powers over many possible distinguishing primes. This answers the low-zero question for quadratic characters in the ProveIt article and extends it to characters of arbitrary order.

\section{The zero data and the results}
For a primitive Dirichlet character \(\chi\) write \(q_\chi\) for its conductor. Include the character of conductor one, for which \(L(s,\chi)=\zeta(s)\). A zero in this paper means one in \(0<\Re s<1\). We use the Dirichlet GRH theorem of \cite{lean} and put
\[
 m_\chi(t)=\ord_{s=1/2+it} L(s,\chi),\qquad
 L=\log(2Q),\quad D(Q)=\frac{L^2}{\log(2L)},\quad
 J(t)=\log\log\log(e^{e^e}+|t|).
\]
The multiplicity is zero at an ordinary point. A sum over \(t\) counts each distinct ordinate once, with the multiplicity supplied by \(m_\chi(t)\). For a complex character the sign of the ordinate has meaning. Only positive ordinates are needed in the low-height conclusion.

\begin{theorem}[Uniform local independence]\label{thm:main}
For every fixed integer \(r\ge1\) there are constants \(C_r,c_r>0\) with the following property. Let \(Q\ge3\), let \(\chi_1,\ldots,\chi_r\) be distinct primitive characters of conductors at most \(Q\), and let \(I=(U-H/2,U+H/2)\). If
\begin{equation}\label{eq:length}
 H\ge C_r\bigl(D(Q)+J(U)\bigr),
\end{equation}
then every \(a=(a_1,\ldots,a_r)\in\C^r\) satisfies
\begin{equation}\label{eq:main}
 \sum_{t\in I}\left|\sum_{j=1}^r a_j m_{\chi_j}(t)\right|
 \ \ge\ c_r H\frac{\log(2L)}{L}\,\|a\|_2 .
\end{equation}
The constants are uniform in the conductors, orders, window centre and coefficients.
\end{theorem}

Taking \(U=H/2\) gives a window \((0,H)\) with \(H=C'_r D(Q)\). Indeed \(J(H/2)=o(H)\), and we may increase the constant to cover all \(Q\ge3\). For two distinct characters at this cutoff the total multiplicity discrepancy is \(\gg\log(2Q)\). It can include a difference of orders at a shared zero. A nonreal character and its conjugate are covered by the same bound.

\begin{corollary}[Recovering products and fields]\label{cor:field}
Fix \(d\ge1\). In a window satisfying \eqref{eq:length} with a constant depending only on \(d\), the zero multiset determines each product
\[
 \mathcal L(s)=\prod_{j=1}^{k} L(s,\chi_j)^{e_j},
 \qquad k\le d,\quad e_j\in\mathbb Z_{>0},\quad q_{\chi_j}\le Q,
\]
within this class, including every character and exponent. Consequently the same data for \(\zeta_K\) determine an abelian field \(K/\Q\) of degree at most \(d\) and conductor at most \(Q\), within that class.
\end{corollary}
\begin{proof}
If two products agree in their local zero data, subtract their vectors of factor exponents on the union of the character sets. Its size is at most \(2d\). Orders add in a product, so \eqref{eq:main} forces this difference to vanish. Use the maximum of the finitely many threshold constants up to \(2d\).

For an abelian field, Kronecker--Weber and abelian Artin factorization give \(\zeta_K=\prod_{\chi\in X_K}L(s,\chi)\), where \(X_K\) has \([K:\Q]\) primitive constituents with conductors dividing the conductor of \(K\) \cite{kw}. If two candidate fields have the same constituents, embed them in the cyclotomic field of level the least common multiple of their conductors. Inflating these characters to that level, the intersection of their kernels fixes precisely \(K\). Thus the local data determine the field as a subfield of \(\Q^{\rm ab}\). The same cutoff applies with an upper bound on the absolute discriminant in place of the conductor bound, since the conductor--discriminant formula gives \(|\operatorname{disc} K|=\prod_{\chi\in X_K}q_\chi\) \cite[Thm.~3.11]{wash}.
\end{proof}

For example, \(\Q(\zeta_5)\) and \(\Q(\sqrt5,\sqrt{-3})\) share the zeta factor and the quadratic character of discriminant 5. The first field's remaining two constituents are the conjugate quartic characters modulo 5. The other field has the quadratic constituents of discriminants \(-3\) and \(-15\). The corollary distinguishes their zeta zero lists even though zeros from the two common factors appear in both lists.

\section{Arithmetic input and the choice of a prime}
Use the Fourier convention
\[
 h(t)=\int_{\R}g(u)e^{itu}\dd u,\qquad
 g(u)=\frac1{2\pi}\int_{\R}h(t)e^{-itu}\dd t.
\]
For a primitive \(\chi\), put \(\kappa\in\{0,1\}\) with \(\chi(-1)=(-1)^\kappa\) and \(\delta_\chi=\one_{q_\chi=1}\). The Guinand--Weil explicit formula in this convention is \cite[eqs.~(2.1), (2.6)]{weil}
\begin{align}
 \sum_\gamma h(\gamma)
 &=\delta_\chi\{h(i/2)+h(-i/2)\}+g(0)\log(q_\chi/\pi)+\mathcal A_\kappa(h)\notag\\
 &\quad-\sum_{n\ge2}\frac{\Lambda(n)}{\sqrt n}
       \{\chi(n)g(\log n)+\overline{\chi(n)}g(-\log n)\}, \label{eq:weil}\\
 \mathcal A_\kappa(h)&=\frac1{2\pi}\int_\R h(t)\Re\frac{\Gamma'}{\Gamma}
             \left(\frac{1/2+\kappa+it}{2}\right)\dd t .\notag
\end{align}
Here and below the sum over \(\gamma\) repeats multiplicities. Analytic tests with sufficient decay in a strip wider than \(|\Im t|\le1/2\) are admissible. The complex form follows by linearity, or directly by using the conjugate character on the reflected contour. When \(g(0)=0\), the integral representation of the digamma function gives
\begin{equation}\label{eq:arch}
 \mathcal A_\kappa(h)=-\int_0^\infty
 \frac{e^{-(1/2+\kappa)u}}{1-e^{-2u}}\,[g(u)+g(-u)]\dd u .
\end{equation}
Indeed the constant terms integrate to multiples of \(g(0)\), and the cosine in the real part inverts to \((g(u)+g(-u))/2\) after the change of variable by a factor of two. We also use the classical unit-interval count
\begin{equation}\label{eq:count}
 \#\{\gamma:|\gamma-v|\le1\}\ll\log(2q_\chi(3+|v|)).
\end{equation}
For conductor greater than one this follows from the Dirichlet zero-count formula with logarithmic remainder, bounding either signed interval by the corresponding interval of absolute heights \cite{count}. The conductor-one case follows from the Riemann--von Mangoldt formula for zeta. The bounded-height range follows from the total count up to height 3.

\begin{lemma}[Many distinguishing primes]\label{lem:many}
Fix \(r\) and a unit vector \(a\in\C^r\). For a sufficiently large constant \(A_r\), set \(X=A_rL^2\) and \(S(n)=\sum_j a_j\chi_j(n)\). At least \(b_r X/\log X\) primes \(p\in[X,2X]\) satisfy \(|S(p)|\ge1/2\), where \(b_r>0\) depends only on \(r\).
\end{lemma}
\begin{proof}
Fix nonnegative \(w\in C_c^\infty((1,2))\) of integral one. Substituting \(g(u)=e^{u/2}w(e^u/X)\) in \eqref{eq:weil} gives for a primitive \(\xi\)
\[
 \sum_n\Lambda(n)\xi(n)w(n/X)=\delta_\xi X+O_w(\sqrt X\log(2q_\xi)).
\]
In detail, \(h(t)=\sqrt X X^{it}\int_1^2 w(y)y^{-1/2+it}\dd y\), whose absolute value on the real axis is \(\ll_w\sqrt X(1+|t|)^{-3}\) by integration by parts. Equation~\eqref{eq:count} bounds the zero sum, Stirling's formula bounds the gamma integral by \(O_w(\sqrt X)\), and the two pole terms are \(X\) and \(O_w(1)\). The real-ordinate estimate is where GRH saves a power of \(X\).

Apply this to the primitive inducer of \(\chi_j\overline{\chi_k}\) on the least common multiple of the two conductors, at most \(Q^2\). It is principal exactly when \(j=k\), by uniqueness of primitive induction. Removing proper prime powers costs \(O(\sqrt X\log^2(2X))\), since their number in \([X,2X]\) is \(O(\sqrt X\log(2X))\). At primes the omitted Euler factors cost at most a constant times \(\log(q_{\chi_j}q_{\chi_k})\). Hence
\[
 \sum_p \log p\,w(p/X)\chi_j(p)\overline{\chi_k(p)}
 =\one_{j=k}X+O_w\big(\sqrt X[L+\log^2(2X)]+L\big).
\]
Choose \(A_r\) large enough that each error is at most \(X/(8r)\), uniformly for \(Q\ge3\). Since \((\sum_j|a_j|)^2\le r\), the weighted sum of \(|S(p)|^2\) is at least \(7X/8\). The same prime estimate for the trivial character gives total weight at most \(5X/4\). The primes with \(|S(p)|<1/2\) therefore contribute at most \(5X/16\). Each remaining prime contributes at most \(r\|w\|_\infty\log(2X)\). This proves the assertion. Increasing \(A_r\) preserves a choice of \(b_r\) independent of that increase.
\end{proof}

\begin{lemma}[Average contribution of neighbouring prime powers]\label{lem:sieve}
For \(p\in[X,2X]\) prime and \(B\ge KX/\log X\), \(K\ge1\), put
\[
 E_p=\sum_{\substack{X/2\le n\le4X\\n\ne p}}
 \frac{\Lambda(n)}{\log X}\min\left(1,\left(\frac{16X}{B|n-p|}\right)^4\right).
\]
With absolute implicit constants,
\begin{equation}\label{eq:average}
 \sum_{X\le p\le2X} E_p
 \ll \frac{X}{K\log X}+\sqrt X\log X(1+\log X).
\end{equation}
\end{lemma}
\begin{proof}
Write \(D=16\log X/K\), which majorizes \(16X/B\), and \(v(h)=\min(1,(D/h)^4)\). For prime neighbours the weight \(\Lambda(n)/\log X\) is bounded. For \(1\le h\le X/4\) the uniform upper-bound sieve for prime pairs bounds the number of pairs with \(|p-n|=h\) by
\[
 \ll \frac{X}{(\log X)^2}\mathfrak S(h),\qquad
 \sum_{1\le h\le y}\mathfrak S(h)\ll y.
\]
Here odd shifts give no pairs for large \(X\). One may take \(\mathfrak S(h)=\prod_{\ell\mid h,\,\ell>2}(\ell-1)/(\ell-2)\) for even \(h\); factors independent of \(h\) are absorbed in the constant. Apply the sieve to the smaller prime on an interval covering \([X/2,2X]\). See \cite[Lemma 2.1]{evans}. Partial summation gives \(\sum_{h\ge1}\mathfrak S(h)v(h)\ll D\), also for \(D<1\). Shifts larger than \(X/4\) contribute \(O(X^2(D/X)^4)\) even if all integers in both ranges are counted. This is absorbed by the second term of \eqref{eq:average}. There are \(O(\sqrt X\log X)\) proper prime powers in the local range. For each, summing over every possible integer \(p\) costs \(O(1+D)\) by comparison with the integral of \(v\). This gives \eqref{eq:average}.
\end{proof}

Given any fixed \(\eta_r>0\), first choose \(K=K_r\) sufficiently large and then \(A_r\) sufficiently large. Lemmas~\ref{lem:many} and \ref{lem:sieve} guarantee a prime with both
\begin{equation}\label{eq:witness}
 |S(p)|\ge\tfrac12,\qquad E_p\le\eta_r.
\end{equation}
Indeed fewer than \(b_r X/(2\log X)\) primes can have \(E_p>\eta_r\). This selection works for arbitrarily high-order characters, whose distinct nonzero values at a particular prime can be arbitrarily close.

\section{A test for one Euler coefficient in a translated window}
Let \(W\) be the fourfold convolution of the uniform probability density on \([-1/4,1/4]\), and let \(k(t)=c\exp(-\cosh t)\), normalized to integral one. Put \(W_B(t)=B^{-1}W(t/B)\), \(\Phi_B=W_B*k\), and here write \(\widehat f(v)=\int f(t)e^{-itv}\dd t\). Then
\[
 F_B(v)=\widehat\Phi_B(v)=\sinc(Bv/4)^4\widehat k(v),\qquad
 \sinc z=\sin(z)/z,\quad \sinc 0=1.
\]
Both densities are even. The support of \(W_B\) is \([-B,B]\). Moving the contour for \(\widehat k\) to imaginary part \(-\operatorname{sgn}(v)\) is valid since \(\cos(1)>0\). It gives
\begin{equation}\label{eq:F}
 |F_B(v)|\le C e^{-|v|}\min(1,(4/(B|v|))^4),\qquad F_B(0)=1.
\end{equation}
For local estimates one may also use \(|\widehat k(v)|\le1\).

For a prime \(p\) satisfying \eqref{eq:witness} let \(\lambda=\log p\), \(\delta=\pi/(2\lambda)\). Take \(X\) large enough that \(\delta\le1\). Define
\begin{align}
 g(u)&=\sin(\delta u)\,e^{-iU(u-\lambda)} F_B(u-\lambda),\label{eq:g}\\
 h(t)&=\pi e^{it\lambda}\{\Phi_B(t-U+\delta)+\Phi_B(t-U-\delta)\}.
 \label{eq:h}
\end{align}
Indeed \(\sin(\delta(\lambda+v))=\cos(\delta v)\), which explains the two translates. These tests satisfy \(g(0)=0\) and \(g(\lambda)=1\). The first identity cancels the conductor term exactly, including when the coefficients \(a_j\) have nonzero sum. The second preserves the chosen prime. In the strip \(|\Im t|\le1\) the kernel \(k(t)\) decays double exponentially as \(|\Re t|\to\infty\). This verifies admissibility in \eqref{eq:weil} for each choice of parameters.
For real \(t\), convolution with \(k,k'\) and \(\|W_B\|_\infty\ll B^{-1}\) give
\begin{equation}\label{eq:sup}
 \|h\|_\infty\ll B^{-1},\qquad
 \|h'\|_\infty\ll (1+\log p)/B.
\end{equation}

\begin{lemma}[Estimates for the test]\label{lem:test}
For \(X\le p\le2X\), \(B\ge 4\), the following bounds are absolute and uniform in \(U\).
\begin{align}
 \sum_{n\ne p,\,n\ge2}\frac{\Lambda(n)}{\sqrt n}|g(\log n)|
 &\ll E_p\frac{\log X}{\sqrt X}+B^{-4}\sqrt X\log(2X),\label{eq:leak}\\
 \sum_{n\ge2}\frac{\Lambda(n)}{\sqrt n}|g(-\log n)|&\ll p^{-1},\label{eq:mirror}\\
 |\mathcal A_\kappa(h)|&\ll p^{-1/2}/B+p^{-1},\qquad
 |h(i/2)|+|h(-i/2)|\ll \sqrt p/B.\label{eq:other}
\end{align}
The part of the zero sum outside \(I=(U-2B,U+2B)\), including endpoints, has absolute mass at most
\begin{equation}\label{eq:tail}
 \sum_{|\gamma-U|\ge2B}|h(\gamma)|
 \ll \log(2q_\chi(3+|U|+B))\exp(-c_0 e^B)
\end{equation}
with an absolute \(c_0>0\).
\end{lemma}
\begin{proof}
For \(n\in[X/2,4X]\), the mean value theorem gives \(|\log(n/p)|\ge |n-p|/(4X)\). Apply \eqref{eq:F} and the definition of \(E_p\). Outside this range \(|\log(n/p)|\ge\log2\), so the bound in \eqref{eq:F} is \(\ll B^{-4}\min(n/p,p/n)\). The sum of \(\Lambda(n)n^{-1/2}\min(n/p,p/n)\) is \(O(\sqrt p\log(2p))\), directly by splitting at \(p\) and using \(\Lambda(n)\le\log n\). This proves \eqref{eq:leak}. For the reflected term, the exponential in \eqref{eq:F} supplies \((np)^{-1}\), and \(\sum\Lambda(n)n^{-3/2}<\infty\).

For \(u>0\), \(|\sin(\delta u)|/(1-e^{-2u})\ll1\) since \(\delta\le1\). In \eqref{eq:arch} the negative argument thus costs \(O(1/p)\). For the positive argument set \(v=u-\lambda\). Then \(e^{-u/2}=p^{-1/2}e^{-v/2}\), and
\[
 \int_{\R} e^{-v/2-|v|}\min(1,(4/(B|v|))^4)\dd v\ll B^{-1}.
\]
For the poles, \(\sup_x|\Phi_B(x\pm i/2)|\le \|W_B\|_\infty
\int_\R |k(t\pm i/2)|\dd t\ll1/B\), whereas the modulation in \eqref{eq:h} costs at most \(\sqrt p\).

Finally if \(2B+j\le |t-U|<2B+j+1\) for a nonnegative integer \(j\), the real arguments of \(k\) in \eqref{eq:h} have absolute value at least \(B+j-1\). Positivity and unit mass of \(W_B\) give \(|h(t)|\ll\exp(-c e^{B+j})\). Apply \eqref{eq:count} on the two unit intervals, then sum over \(j\). The extra growth of the logarithm is at most \(O(\log(2+j))\), absorbed by the double exponential. This proves \eqref{eq:tail}.
\end{proof}

\begin{proof}[Proof of Theorem~\ref{thm:main}]
By homogeneity suppose \(\|a\|_2=1\). Set \(B=H/4\) and choose \(X=A_r L^2\). In the sum of \eqref{eq:weil} with coefficients \(a_j\), the main term at \(p\) is
\[
 -\frac{\log p}{\sqrt p} S(p).
\]
Its modulus is at least \((2\sqrt2)^{-1}\log X/\sqrt X\). Each other term in Lemma~\ref{lem:test} can be summed with a factor at most \(\|a\|_1\le\sqrt r\). In the reflected prime channel the coefficients are \(\sum_j a_j\overline{\chi_j(n)}\), also bounded by \(\sqrt r\). Take \(\eta_r\) small for the first term of \eqref{eq:leak}, then \(K_r\) large in \(B\ge K_r X/\log X\). In particular \(\sqrt p/B\ll K_r^{-1}\log X/\sqrt X\), so the possible poles of the conductor-one factor fit in the same error budget. Increasing \(A_r\) handles the terms \(O_r(p^{-1})\), \(O_r(B^{-4}\sqrt X\log(2X))\) and \(O_r(p^{-1/2}/B)\). With these choices their total, including the pole and local contributions, is at most a quarter of the lower bound on the main term.

We detail the height term because uniformity in \(U\) matters here. After fixing \(A_r,K_r\), \eqref{eq:length} with large enough \(C_r\) ensures both \(B\ge K_r X/\log X\) and \(B\ge 2J(U)+C_r'\). In particular
\[
 \log\log(e^{e^e}+|U|)=e^{J(U)}\ll e^{B/2}.
\]
The logarithm of the prefactor in \eqref{eq:tail} is
\(O(1+\log\log(2Q)+\log\log(3+B)+e^{B/2})\);
this is \(o(e^B)\) uniformly in the permitted range as \(B\to\infty\). Thus the summed tail in \eqref{eq:tail} is smaller than a further quarter of the main term, increasing the absolute lower threshold for \(B\) if needed. There is no condition on the ordinates at the endpoints. We obtain
\begin{equation}\label{eq:observable}
 \left|\sum_{t\in I} h(t)\sum_j a_jm_{\chi_j}(t)\right|
 \gg_r \frac{\log X}{\sqrt X}.
\end{equation}
Use \eqref{eq:sup} and \(B=H/4\). With fixed \(A_r\), \(\log X/\sqrt X\asymp_r\log(2L)/L\), giving \eqref{eq:main}.
\end{proof}

\section{Stability under missing and displaced zeros}
The supremum and derivative bounds for the same test give a quantitative interpretation when a measurement lists a nearby ordinate in place of the actual one.

\begin{corollary}[Partial matching]\label{cor:matching}
In a window of Theorem~\ref{thm:main} for two distinct characters, take any partial one-to-one matching of their lists, retaining repeated copies. Let \(M\) count all unmatched entries in both lists. Then
\begin{equation}\label{eq:matching}
 M+\sum_{(t,t')\ {\rm matched}}
 \min\{1,\log(2L)|t-t'|\}\ \gg\ H\frac{\log(2L)}{L}.
\end{equation}
In particular every perfect matching has total displacement \(\gg H/L\).
\end{corollary}
\begin{proof}
Use \eqref{eq:observable} with \((a_1,a_2)=(1,-1)/\sqrt2\). An unmatched entry contributes \(O(1/B)\), and a matched pair contributes at most
\(\min(2\|h\|_\infty,\|h'\|_\infty |t-t'|)\). Since \(\log p\asymp\log(2L)\), use \eqref{eq:sup} and multiply by \(B\). The final assertion follows by dropping the cap for a perfect matching.
\end{proof}

At the low cutoff the combined cost on the left of \eqref{eq:matching} is \(\gg\log(2Q)\). Exact matches have zero cost. An insertion or deletion costs one, and a small displacement costs in proportion to its size. Exact identification refers to the collection of all possible characters or fields under the stated bounds. The constants here have been left unevaluated. A practical decoding radius at a particular conductor would require explicit constants and certified zero data.

\section{Comparison with the sources}
\textbf{The analytic input.} The readable edition of \cite{lean} exposes \texttt{LFunction\_eq\_zero\_re} for every Dirichlet character of positive modulus in the open critical strip, its primitive specialization, and Mathlib's \texttt{RiemannHypothesis}. These are precisely the zero-location inputs for \eqref{eq:weil} and the primitive quotient characters in Lemma~\ref{lem:many}. Andrew Sutherland records independent builds, kernel replay of both editions, and successful comparator checks in \cite{audit}. This note uses that verification record. Theorem~\ref{thm:main} and its corollaries are proved on paper here. The published machine checks cover the upstream input.

\textbf{What the Article asks.} The ProveIt README describes a 621-page synthesis of 103 reports and records the limits of its internal checking \cite{article}. The section on recovering zeros and characters proves the quadratic cutoff \(O((\log Q)^2\log\log Q)\) in Corollary \texttt{rec:cor:quadsep}. Its open question asks whether that last factor can be removed and what height distinguishes higher-order characters and their conjugates. Its algorithm recovers every relevant small prime coefficient by isolating integer logarithms. Lemma~\ref{lem:sieve} uses the sparsity of prime pairs to select from many coefficients. It yields a saving of order \((\log\log Q)^2\) in the sufficient height for identification. The Article also gives triple-logarithmic windows for fixed families (Theorem \texttt{zer:thm:short} and Theorem \texttt{rec:thm:localrank}). Those results include different and sometimes stronger fixed-family conclusions. The contribution here is a common window uniform in conductor and character order, reaching down to zero, with quantitative complex rank and recovery of whole products.

\textbf{Earlier literature.} Srinivas \cite{srinivas} proves a large-height multiplicity separation theorem for fixed functions in the Selberg class. Banks \cite{banks} gives an unconditional, modulus-uniform result for two primitive characters of the same modulus \(q\), in windows of order \(q^\theta\log T\) at \(T\gg_\theta q^\theta\), for \(\theta>1/3\), or \(\theta>1/4\) for cubefree \(q\). Eid's July 2026 preprint \cite{eid} already describes stable polylogarithmic-height distinction for characters of bounded order. Corollary 1.2 of the revised paper gives \(O_{d,\varepsilon}((\log(3Q))^{2+\varepsilon})\) for orders dividing a fixed \(d\), using the signed multiset in \([-T,T]\). Banks and Loftus \cite{loftus} study genuinely unshared zeros of fixed functions at large heights. Our lower bound uses multiplicities, including excess at a shared zero. This comparison identifies the specific extension proved here. Priority beyond the sources reviewed here remains open.

\section{Two reproducible illustrations}
\textbf{Signed ordinates.} Let \(\chi\) be the quartic character modulo 5 with \(\chi(2)=i\). For this example only, use the compact arithmetic test
\[
 g(\log 2+v)=\one_{|v|\le1/4}\cos^4(2\pi v)\sinc(3v)^4.
\]
It sees the prime 2 alone. The conductor and archimedean terms agree for \(\chi\) and \(\overline\chi\). Subtraction in \eqref{eq:weil} gives the exact full imaginary part \(-\sqrt2\log2\). The script independently searches for sign changes of the rotated Hardy function, refines them, and checks residuals at the computed ordinates. A mesh scan is a numerical illustration and supplies no completeness certificate.
\begin{center}
\begin{tabular}{rrrr}
\toprule
Cutoff & Negative / positive found & Imaginary difference & Absolute error\\
\midrule
40 & 16 / 15 & $-0.980528847988$ & $2.71\times10^{-4}$\\
60 & 28 / 27 & $-0.980287688455$ & $2.95\times10^{-5}$\\
100 & 53 / 54 & $-0.980257280235$ & $8.63\times10^{-7}$\\
\bottomrule
\end{tabular}
\end{center}
The four ordinates found with \(|\gamma|<10\) are approximately \(-9.442931,-4.132904,6.183578,8.457229\). Conjugation reflects this list. Taking absolute values would lose the orientation that identifies the character.

\textbf{Selecting a prime by averaging.} Compare the quartic character above with the quadratic character modulo 3. At \(X=4000\), the script finds 346 primes in \([X,2X]\) at which their difference has modulus at least one. For each it sums the absolute neighbouring prime-power contributions in \([X/2,4X]\), using the fourth-power sinc factor and the bound \(|\widehat k|\le1\), and divides by the magnitude of the target term. With \(B=32X/\log X\), the median ratio is \(0.00217\), the largest is \(0.781\), and 304 of the 346 primes have ratio at most \(1/4\). This illustrates how widely the interference varies between primes. The code also cross-checks one row by independent integer-by-integer trial division. The remaining terms of the analytic test are bounded in Lemma~\ref{lem:test}.

\section{Further research questions}
\begin{enumerate}
\item Can the sufficient height \(D(Q)\) be improved uniformly for all primitive characters? The selection argument pays for absolute values of the off-diagonal prime coefficients. A further gain may require cancellation among them, or a different way of using the whole prime signal.
\item How do the least valid constants grow with the number of distinct factors? Quantifying this dependence would allow a degree bound for the abelian fields that grows with \(Q\).
\item Can an analogue at this conductor scale detect a zero at which just one of a pair of functions vanishes? The present multiplicity estimate also counts unequal positive orders.
\item What explicit height and measurement precision work at a specified conductor bound? This requires numerical constants in the sieve and smoothed prime estimate as well as a certified enumeration of the observed zeros.
\end{enumerate}

\begin{thebibliography}{99}\small
\bibitem{lean} Repository \texttt{lsernau/grh-lean}, \emph{GRH for Dirichlet L-functions, in Lean 4}, based on the proof of L. Alp\"oge with Claude. Commit \texttt{72fead102f66}. \href{https://github.com/lsernau/grh-lean/blob/72fead102f663d5c7a407964a13dc0cf9d977207/grh/src/Statements.lean}{Statements} and \href{https://github.com/lsernau/grh-lean/blob/72fead102f663d5c7a407964a13dc0cf9d977207/grh/src/Main.lean}{exported theorems}.
\bibitem{audit} A. V. Sutherland, \emph{Verification files for the shorter Lean proof of RH and of GRH for Dirichlet L-functions}, November 2026. \url{https://math.mit.edu/~drew/shortgrh/}.
\bibitem{article} V. Reshetnikov, ProveIt repository, \emph{Beyond the Riemann hypothesis for Dirichlet L-functions: consequences, transfers and limits}. Snapshot \texttt{dec3e4f21630}. \href{https://github.com/VladimirReshetnikov/ProveIt/tree/dec3e4f216305071186a1cc4c42fe1cdb38df9be/Analysis/GeneralizedRiemannHypothesis/Article}{Article}, \href{https://github.com/VladimirReshetnikov/ProveIt/blob/dec3e4f216305071186a1cc4c42fe1cdb38df9be/Analysis/GeneralizedRiemannHypothesis/Article/sections/rec.tex}{reconstruction section}, and \href{https://github.com/VladimirReshetnikov/ProveIt/blob/dec3e4f216305071186a1cc4c42fe1cdb38df9be/Analysis/GeneralizedRiemannHypothesis/Article/sections/zer.tex}{zero-distribution section}.
\bibitem{weil} A. Chirre, V. J. Pereira J\'unior and D. de Laat, Primes in arithmetic progressions and semidefinite programming, \emph{Math. Comp.} \textbf{90} (2021), 2235--2246. \href{https://doi.org/10.1090/mcom/3638}{DOI}, \href{https://arxiv.org/abs/2005.02393}{preprint}.
\bibitem{count} M. A. Bennett, G. Martin, K. O'Bryant and A. Rechnitzer, Counting zeros of Dirichlet \(L\)-functions, \emph{Math. Comp.} \textbf{90} (2021), 1455--1482. \href{https://doi.org/10.1090/mcom/3599}{DOI}.
\bibitem{evans} N. Evans, Correlations of almost primes, \emph{Math. Proc. Cambridge Philos. Soc.} \textbf{174} (2023), 301--344, Lemma 2.1. \href{https://doi.org/10.1017/S0305004122000251}{DOI}. See also H. L. Montgomery and R. C. Vaughan, \emph{Multiplicative Number Theory I}, Cambridge, 2007, Corollary 3.14.
\bibitem{kw} A. V. Sutherland, \emph{Number Theory I}, MIT 18.785, 2021, lectures \href{https://ocw.mit.edu/courses/18-785-number-theory-i-fall-2021/mit18_785f21_lec19.pdf}{19} and \href{https://ocw.mit.edu/courses/18-785-number-theory-i-fall-2021/mit18_785f21_lec20.pdf}{20}.
\bibitem{wash} L. C. Washington, \emph{Introduction to Cyclotomic Fields}, 2nd ed., GTM 83, Springer, 1997.
\bibitem{srinivas} K. Srinivas, Distinct zeros of functions in the Selberg class, \emph{Acta Arith.} \textbf{103} (2002), 201--207. \href{https://doi.org/10.4064/aa103-3-1}{DOI}.
\bibitem{banks} W. D. Banks, Comparing zeros of distinct Dirichlet \(L\)-functions, \emph{Publ. Math. Debrecen} \textbf{106} (2025), 327--335. \href{https://publi.math.unideb.hu/paper/4308}{Published article}.
\bibitem{eid} S. Eid, \emph{Effective Distinguishing of Primitive Dirichlet Characters by Their Low-Lying Zeros}, revised PDF, July 2026. Deposited as \emph{Effective Determination of Primitive Dirichlet Characters by Their Zeros up to Polylogarithmic Height}. \href{https://doi.org/10.5281/zenodo.21326653}{Version DOI}.
\bibitem{loftus} W. D. Banks and K. Loftus, \emph{Unshared zeros of Dirichlet L-functions}, 2026. \href{https://arxiv.org/abs/2607.22930}{arXiv:2607.22930}.
\end{thebibliography}
\end{document}
